\section*{Chapter \ref{ch:dv:fourier-pricing-and-calibration-engineering} --- Fourier Pricing and Calibration Engineering}

\begin{solution}{exo:dv:fourier-pricing-and-calibration-engineering:1}
$|\varphi(u-\frac i2)|\le\E[e^{x_T/2}]\le1$ by Jensen's inequality, since $\E[e^{x_T}]=1$. The integral is therefore bounded by
$\int_0^\infty du/(u^2+\frac14)=\pi$, and $\sqrt{FK}\to0$ gives $C\to P(0,T)F$. $\varphi(-i)=\E[e^{x_T}]=\E[S_T]/F=1$ says that the forward is the
mean of $S_T$ under the pricing measure, as it must be for a forward-martingale price.
\end{solution}

\begin{solution}{exo:dv:fourier-pricing-and-calibration-engineering:2}
As $k\to-\infty$ the call tends to $F$, so it is not integrable in $k$ and has no Fourier transform. The factor $e^{\alpha k}$ makes it decay on
that side. The transform then involves $\varphi(v-(\alpha+1)i)$, that is $\E[(S_T/F)^{\alpha+1}]$, which must be finite: $\alpha$ is limited by
the moments the model has. In Heston with a strongly negative correlation and a long expiry, high moments explode in finite time.
\end{solution}

\begin{solution}{exo:dv:fourier-pricing-and-calibration-engineering:3}
The model price is Black's formula at the model's implied volatility: $C^{\text{model}}=C^{\text{Black}}(\sigma^{\text{model}})$. Expanding around
$\sigma^{\text{mid}}$, $C^{\text{model}}-C^{\text{mid}}=\mathcal V(\sigma^{\text{model}}-\sigma^{\text{mid}})+\frac12\,\partial_\sigma\mathcal
V\,(\sigma^{\text{model}}-\sigma^{\text{mid}})^2+\dots$ Dividing by $\mathcal V$ leaves the volatility error plus a second-order term (the \omterm{def:dv:greeks-and-the-hedging-pnl:greeks}{volga}).
\end{solution}

\begin{solution}{exo:dv:fourier-pricing-and-calibration-engineering:4}
With $n$ quotes of independent noise $\sigma$ and $p$ fitted parameters, a correct model leaves residuals of root-mean-square
$\sigma\sqrt{(n-p)/n}$: $0.3\sqrt{13/18}=0.255$ volatility point for eighteen quotes after the filter. The chapter's fits give 0.22 on the first
day and 0.291 on average: noise plus a little misspecification, since the market is a \omterm{def:dv:jumps-and-levy-models:bates}{Bates model}. A fit well below the noise floor would be
fitting the noise.
\end{solution}

\begin{solution}{exo:dv:fourier-pricing-and-calibration-engineering:5}
COS prices the option under the density truncated to $[a,b]$. The mass outside, in the heavy left tail of a \omterm{def:dv:stochastic-volatility:heston}{Heston model} with $\rho=-0.7$ and
$\eta=0.6$, is simply lost. More terms converge to the price under the truncated density, not the true one. Only a wider interval (sixteen
standard deviations) removes the error, at the cost of more terms for the same resolution.
\end{solution}

\begin{solution}{exo:dv:fourier-pricing-and-calibration-engineering:6}
The \omterm{def:dv:stochastic-volatility:sv}{volatility of volatility} shows up in the smile's curvature and skew, damped by mean reversion. Over a horizon $T$, the variance's response to
its own shocks is averaged by the factor $(1-e^{-\kappa T})/(\kappa T)$. A larger $\eta$ with a faster $\kappa$ gives about the same effective
volatility of variance at the quoted expiries (up to a year). The quotes see the product, not the two separately, hence the valley.
\end{solution}

\begin{solution}{exo:dv:fourier-pricing-and-calibration-engineering:7}
With equal price weights the largest daily move of $\eta$ is 0.255 (0.235 with \omterm{def:dv:greeks-and-the-hedging-pnl:greeks}{vega} weights), and the standard deviation of the daily moves 0.066
(0.089). The wander is about the same. The weights change which quotes pin the parameters (with equal weights, the one-year options, whose \omterm{def:dv:greeks-and-the-hedging-pnl:greeks}{vegas}
are largest), not the flatness of the valley. The penalty is still needed.
\end{solution}

\begin{solution}{exo:dv:fourier-pricing-and-calibration-engineering:8}
The fit error measures only the vanillas. In the chapter the unpenalised fits stay below 0.45 volatility point every day (at most 0.42), while
$\eta$ moves by 0.235 in a day and a \omterm{def:dv:variance-swaps-and-volatility-derivatives:volswap}{volatility swap}'s convexity adjustment by 0.93 volatility point. Check the parameter paths, their alarms and
the \omterm{def:dv:fourier-pricing-and-calibration-engineering:recalib}{recalibration P\&L} of the exotics too.
\end{solution}

\begin{solution}{pb:dv:fourier-pricing-and-calibration-engineering:1}
\textbf{1.} COS: it prices the quoted strikes directly, and its error falls exponentially with the terms. The FFT gives a fixed grid and needs
interpolation.
\textbf{2.} 128 terms ($8.6\times10^{-5}$).
\textbf{3.} The cumulants: $[a,b]=c_1+\ln(F/K)\pm L\sqrt{c_2}$ with $L=16$. Too narrow ($L=12$) and the error stalls at $3.0\times10^{-5}$.
\textbf{4.} About eight iterations of six residual evaluations, each pricing four smiles: some two hundred smiles, 800 characteristic-function
evaluations each. Those evaluations (complex exponentials and logarithms) dominate; in the build a calibration takes well under a second.
\textbf{5.} Against the Gauss--Legendre Lewis reference, against Black's formula with a lognormal characteristic function, and by $\varphi(0)=
\varphi(-i)=1$.
\textbf{6.} Two quotes, one crossed and one stale (wide); eighteen remain.
\textbf{7.} $v_0=0.0451$, $\kappa=1.82$, $\bar v=0.0518$, $\eta=0.488$, $\rho=-0.632$; 0.22 volatility point.
\textbf{8.} Its error rises from 0.35 to 0.46 volatility point.
\textbf{9.} Between $\eta=0.35$ and 0.65 the best error stays between 0.23 and 0.29 volatility point. Over the whole range from 0.25 to 1 it varies by
less than the 0.3 point of quote noise.
\textbf{10.} $v_0$, $\rho$ and $\eta$ are identified (standard deviations 0.0007, 0.046 and 0.049); $\kappa$ is not (0.56, 28\% of its value).
\textbf{11.} 0.235, from day 58 to day 59 (0.51 to 0.28).
\textbf{12.} It falls from 1.32 to 0.39 volatility point, a change of 0.93: 93\,000 for a \omterm{def:dv:variance-swaps-and-volatility-derivatives:veganot}{vega notional} of 100\,000.
\textbf{13.} The quotes barely changed (the fit error is 0.31 on both days); the parameters slid along the valley. No risk factor that a hedge covers
moved.
\textbf{14.} The largest daily move of $\eta$ falls to 0.071; the convexity adjustment's daily changes fall from 0.36 to 0.07 volatility point
(standard deviation), and on day 59 it moves by 0.18 instead of 0.93.
\textbf{15.} The average fit error rises from 0.291 to 0.303 volatility point, by 0.012.
\textbf{16.} For each $\lambda$ on a grid from $10^{-5}$ to $10^{-2}$, rerun the sixty days with the production pipeline; record the largest daily move of
$\eta$ and the average fit error; take the smallest $\lambda$ that meets both limits.
\textbf{17.} The parameters lag a genuine change in the market (a real rise in the \omterm{def:dv:stochastic-volatility:sv}{volatility of volatility} after a shock). The fit error rises and the
model misprices until it catches up. The fit-error alarm must be able to override the penalty.
\textbf{18.} Penalise the poorly identified directions ($\kappa$, $\eta$) more, and the well-identified ones ($v_0$, $\rho$) less, so that these follow
the market every day.
\textbf{19.} $\lambda=10^{-3}$: the largest daily move of $\eta$ falls from 0.235 to 0.071 (at $3\times10^{-4}$ it is still 0.133), and the average fit
error rises by 0.012 volatility point.
\textbf{20.} Parameters that move only when the quotes require it, so that exotic prices follow the market and not the noise.
\end{solution}

\begin{solution}{iq:dv:fourier-pricing-and-calibration-engineering:1}
Write the price as an integral in Fourier space: the \omterm{def:dv:fourier-pricing-and-calibration-engineering:lewis}{Lewis formula} (Parseval's identity with the payoff's transform on a strip), the Carr--Madan
damped-call transform evaluated by FFT, or the COS expansion of the density. All need only $\varphi$, and all are checked by $\varphi(0)=\varphi(-i)=1$
and by a lognormal case against Black.
\iqlookfor{one of the three routes, and a check.}
\end{solution}

\begin{solution}{iq:dv:fourier-pricing-and-calibration-engineering:2}
Carr--Madan: a whole log-strike grid in one FFT, but a damping parameter to choose, interpolation to the quoted strikes, and errors from the grid's
spacing and range. COS: any strikes, exponential convergence for smooth densities, few terms; it needs a truncation interval from the cumulants,
and heavy tails need a wide one.
\iqlookfor{grid versus strikes, damping versus truncation, convergence.}
\end{solution}

\begin{solution}{iq:dv:fourier-pricing-and-calibration-engineering:3}
The problem is ill-posed: a valley of near-equal fits (for Heston, $\kappa$ and $\eta$ together), along which noise moves the optimum. Measure it
(profile, refits with noise, the Jacobian's condition number). Then add a penalty on changes from yesterday, chosen by replay, and fix or penalise
the poorly identified parameters; watch the exotics' \omterm{def:dv:fourier-pricing-and-calibration-engineering:recalib}{recalibration P\&L}.
\iqlookfor{ill-posedness, penalty, identification.}
\end{solution}

\begin{solution}{iq:dv:fourier-pricing-and-calibration-engineering:4}
Store the quote snapshot, the filter's decisions, settings, the start point, the result and diagnostics, so that any day can be replayed. Check the
filter counts, the fit error overall and by expiry, the condition number and the parameter moves. Alarm when a parameter moves beyond its limit or
the fit error beyond its own, and hold the previous parameters until a person releases the new ones.
\iqlookfor{replayability, diagnostics, and an alarm that blocks.}
\end{solution}

\begin{solution}{iq:dv:fourier-pricing-and-calibration-engineering:5}
If the model's parameters jumped along a flat direction, the exotics' values moved with no move in any hedged risk factor: \omterm{def:dv:fourier-pricing-and-calibration-engineering:recalib}{recalibration P\&L}, which
the attribution leaves unexplained. Check the parameter history for that day, reprice with yesterday's parameters on today's market, and compare.
\iqlookfor{\omterm{def:dv:fourier-pricing-and-calibration-engineering:recalib}{recalibration P\&L}, and repricing with yesterday's parameters.}
\end{solution}

\begin{solution}{iq:dv:fourier-pricing-and-calibration-engineering:6}
Unweighted price errors let long-dated at-the-money options dominate. Weight by inverse \omterm{def:dv:greeks-and-the-hedging-pnl:greeks}{vega} to fit in volatility, and further by inverse bid--ask
spread (or quote reliability), so that liquid quotes count more. Then add weights for what the book needs, such as the expiries and strikes of its
exotics' hedges.
\iqlookfor{\omterm{def:dv:greeks-and-the-hedging-pnl:greeks}{vega}, spread, and the book's needs.}
\end{solution}
